\chapter{绝对值与完备化}\label{chap:nt-4}
\section{绝对值与数域的位}\label{sec:nt-4-1}
\subsection{定义}\label{sec:nt-4-1-1}
域 $K$ 上的一个\newterm{nt-absolute-value}{绝对值}{absolute value}是映射 $|\cdot|\colon K\longrightarrow\mathbb R_+$，使得对任意 $x,y$ 都有
\begin{enumerate}[label=\roman*)]
\item $|x|=0\Longleftrightarrow x=0$.
\item $|xy|=|x||y|$.
\item $|x+y|\le |x|+|y|$.
\end{enumerate}
有些作者称之为乘法范数，一些英语作者则称之为赋值.这里采用 Bourbaki 的术语.
若绝对值满足更强的公理
\begin{enumerate}[label=\roman*)\textquotesingle,start=3]
\item $|x+y|\le\max(|x|,|y|)$，
\end{enumerate}
则称它为\newterm{nt-nonarchimedean-absolute-value}{非 Archimedes 绝对值}{non-Archimedean absolute value}.
\begin{lemma*}
以下条件等价：
\begin{enumerate}[label=\roman*)]
\item $|\cdot|$ 是非 Archimedes 的.
\item $\forall n\in\mathbb Z$，$|n|\le1$（仍用 $n$ 表示 $n$ 在 $K$ 中的像）.
\item $\exists M>0$，$\forall n\in\mathbb Z$，$|n|\le M$.
\end{enumerate}
\end{lemma*}
\begin{proof}
\textup{i)$\Rightarrow$ii)} 事实上，$|n|=|1+1+\cdots+1|\le\max(1,1,\ldots,1)=1$.
\textup{ii)$\Rightarrow$iii)} 是显然的.
\textup{iii)$\Rightarrow$i)} 有
\begin{align*}
|x+y|^n
&=\left|\sum_{k=0}^n\binom nk x^ky^{n-k}\right|
\le M\sum_{k=0}^n|x^ky^{n-k}|\\
&\le M(n+1)\max(|x|,|y|)^n,
\end{align*}
从而
\begin{align*}
|x+y|\le\lim_{n\to\infty}(M(n+1))^{1/n}\max(|x|,|y|)=\max(|x|,|y|).
\end{align*}
\end{proof}
\emph{推论.}若 $K$ 的特征大于 $0$，则它的所有绝对值都是非 Archimedes 的.
\begin{remark*}
若 $|\cdot|$ 是 $K$ 上的非 Archimedes 绝对值，则：
\begin{itemize}
\item $\mathcal O:=\{x\in K,\ |x|\le1\}$ 是一个环；
\item $\mathfrak m:=\{x\in K,\ |x|<1\}$ 是它唯一的极大理想.
\end{itemize}
\end{remark*}
\emph{例.}
\begin{enumerate}[label=\roman*)]
\item 平凡绝对值：$|x|=1$，$\forall x\in K^\times$.
\item $\mathbb C$ 上的模 $z\mapsto |z|=\sqrt{z\overline z}$.
\item 若 $v\colon K\longrightarrow\mathbb R\cup\{\infty\}$ 是一个赋值，则对任意实数 $\rho>1$，$x\mapsto\rho^{-v(x)}$ 都是非 Archimedes 绝对值.
\end{enumerate}
\emph{数域.}设 $K$ 是数域.
\begin{itemize}
\item 对每个嵌入 $\sigma\colon K\hookrightarrow\mathbb C$，定义相应的Archimedes 绝对值 $|x|_\sigma:=|\sigma(x)|$.
\item 对每个 $\mathfrak p\in\operatorname{Max}(\mathcal O_K)$，定义相应的非 Archimedes 绝对值 $|x|_{\mathfrak p}:=N\mathfrak p^{-v_{\mathfrak p}(x)}$.
\end{itemize}
\subsection{绝对值定义的拓扑}\label{sec:nt-4-1-2}
这是由开球
\begin{align*}
B(x,r):=\{y\in K,\ |x-y|<r\},\qquad x\in K,\quad r>0,
\end{align*}
生成的拓扑，即距离 $d(x,y):=|x-y|$ 所对应的度量拓扑.$K$ 的乘法和加法都是连续的；我们称 $K$ 为拓扑域.
\emph{非 Archimedes 情形的特点.}此时距离 $d$ 是超度量，因此有以下两个看似奇特的性质：
\begin{itemize}
\item 球中的每一点都是球心：$\forall y\in B(x,r)$，$B(x,r)=B(y,r)$.
\item 若两个球 $B_1,B_2$ 的交非空，则 $B_1\subset B_2$ 或 $B_2\subset B_1$.
\end{itemize}
\emph{离散绝对值的特点.}若 $|K^\times|$ 是 $\mathbb R_+^\times$ 的离散子群，则开球也是闭集，半径非零的闭球也是开集.因此，这个拓扑全不连通（连通子集只有单点集）.
\begin{example*}
若 $K$ 是数域且 $\mathfrak p\in\operatorname{Max}(\mathcal O_K)$，则与 $|\cdot|_{\mathfrak p}$ 对应的拓扑称为 $\mathfrak p$-进拓扑.一组开集基由
\begin{align*}
x+\mathfrak p^i=B\bigl(x,(N\mathfrak p)^{-i}\bigr),\qquad x\in K,\quad i\in\mathbb N,
\end{align*}
给出.\footnote{译注：原文此处的球定义使用严格不等号，而所列等式的右端不含满足 $|y-x|_{\mathfrak p}=(N\mathfrak p)^{-i}$ 的边界点；此外，$\mathfrak p^i$ 作为 $\mathcal O_K$ 的理想与局部化中的理想须作区分.这里保留原式.}
\end{example*}
\begin{lemma*}
对两个非平凡绝对值 $|\cdot|_1$ 与 $|\cdot|_2$，以下条件等价：
\begin{enumerate}[label=\roman*)]
\item $|\cdot|_1$ 与 $|\cdot|_2$ 定义相同的拓扑.
\item $\forall x\in K$，$|x|_1\le1\Rightarrow |x|_2\le1$.
\item $\exists a>0$，$|\cdot|_2=|\cdot|_1^a$.
\end{enumerate}
\end{lemma*}
\begin{proof}
\textup{i)$\Rightarrow$ii)} 由 $|x|<1\Longleftrightarrow\lim_{n\to\infty}x^n=0$ 即可看出.
\textup{ii)$\Rightarrow$iii)} 由于 $|\cdot|_1$ 非平凡，可取 $y\in K$ 使 $|y|_1>1$.于是 $|y|_2>1$，且存在唯一的 $a>0$ 使 $|y|_2=|y|_1^a$.现在任取 $x\in K$，并设 $b$ 满足 $|x|_1=|y|_1^b$.需要证明 $|x|_2=|y|_2^b$.设 $r=\frac mn\ge b$.则 $|x|_1\le|y|_1^r$，所以 $|x^n/y^m|_1\le1$，故 $|x^n/y^m|_2\le1$，最终得到 $|x|_2\le|y|_2^r$.令 $r$ 趋于 $b$，得到 $|x|_2\le|y|_2^b$.同理可证 $|x|_2\ge|y|_2^b$.
\textup{iii)$\Rightarrow$i)} 在 iii) 的条件下，这两个绝对值定义相同的开球.
\end{proof}
\begin{definition*}
若引理中的条件成立，则称 $|\cdot|_1$ 与 $|\cdot|_2$ 等价.$K$ 上非平凡绝对值的一个等价类称为 $K$ 的位.
\end{definition*}
\subsection{Ostrowski 定理}\label{sec:nt-4-1-3}
\begin{theorem*}[Ostrowski]
$\mathbb Q$ 上的非平凡绝对值都等价于通常的绝对值 $|\cdot|_\infty$，或者等价于某个素数 $p$ 对应的 $p$-进绝对值 $|\cdot|_p$.
\end{theorem*}
\begin{proof}
设 $|\cdot|$ 是 $\mathbb Q$ 上的非平凡绝对值，且 $(a,b)\in\mathbb N$、$a>1$.将 $b$ 按 $a$ 进制展开：
\begin{align*}
\exists!m\in\mathbb N\text{ 及 }\exists!(b_0,\ldots,b_m)\in\{0,\ldots,a-1\}^m,
\qquad b=b_0+b_1a+\cdots+b_ma^m.
\end{align*}
注意 $m\le\frac{\log(b)}{\log(a)}$.令 $M_a:=\max\{|1|,|2|,\ldots,|a-1|\}$，可得
\begin{align*}
|b|&\le\sum_{i=0}^m|b_i||a|^i
\le(m+1)M_a\max(1,|a|^m)\\
&\le M_a\left(\frac{\log(b)}{\log(a)}+1\right)
\max\left(1,|a|^{\frac{\log(b)}{\log(a)}}\right).
\end{align*}
将这个不等式用于 $b^n$，得到
\begin{align*}
|b|&\le\lim_{n\to\infty}M_a^{1/n}
\left(\frac{n\log(b)}{\log(a)}+1\right)^{1/n}
\max\left(1,|a|^{\frac{\log(b)}{\log(a)}}\right)\\
&=\max\left(1,|a|^{\frac{\log(b)}{\log(a)}}\right).
\end{align*}
于是分为两种情形：
\begin{enumerate}[label=\roman*)]
\item 存在 $b\in\mathbb N$ 使 $|b|>1$.此时 $|a|>1$，交换两者的角色，得到
\begin{align*}
|b|^{1/\log(b)}=|a|^{1/\log(a)}=:\rho.
\end{align*}
于是，对任意 $x>0$ 都有 $|x|=\rho^{\log(x)}$，从而对任意 $x\in\mathbb Q$ 都有 $|x|=|x|_\infty^{\log(\rho)}$.
\item 对任意 $b\in\mathbb N$，都有 $|b|\le1$.此时，环 $\mathcal O=\{b\in\mathbb Q,\ |b|\le1\}$ 是包含 $\mathbb Z$ 的局部环.设 $\mathfrak m=\{b\in\mathbb Q,\ |b|<1\}$ 为其极大理想.则 $\mathbb Z\cap\mathfrak m$ 是素理想，且因为 $|\cdot|$ 非平凡，它非零.所以它形如 $(p)$，其中 $p$ 为素数.容易看出，对任意 $x\in\mathbb Q$，$|x|_p\le1\Rightarrow |x|\le1$.由前一引理的判别条件，得到 $|\cdot|$ 等价于 $|\cdot|_p$.
\end{enumerate}
\end{proof}
\begin{remark*}
用 $|p|_p=p^{-1}$ 来规范化 $|\cdot|_p$ 的好处在于，可以得到\newterm{nt-product-formula}{乘积公式}{product formula}（以 $\mathcal P$ 表示素数集合）：\footnote{译注：原文的量词包括 $x=0$；乘积公式通常应限定 $x\ne0$.}
\begin{align*}
\forall x\in\mathbb Q,\qquad\prod_{v\in\mathcal P\cup\{\infty\}}|x|_v=1.
\end{align*}
\end{remark*}
\subsection{完备化：存在性与性质}\label{sec:nt-4-1-4}
我们知道，$\mathbb Q$ 对 $|\cdot|_\infty$ 并不完备，而 $\mathbb R$ 是它的完备化.同样，我们将会看到 $\mathbb Q$ 对 $p$-进拓扑并不完备，因而值得将其完备化.在此之前，先给出完备化的定义与构造.
\emph{赋值域.}若 $K$ 是一个域，$|\cdot|$ 是其上的绝对值，则称二元组 $(K,|\cdot|)$ 为\newterm{nt-valued-field}{赋值域}{valued field}.赋值域的态射
\begin{align*}
(K,|\cdot|_K)\longrightarrow(L,|\cdot|_L)
\end{align*}
是域态射 $\sigma$，使对所有 $x\in K$ 都有 $|\sigma(x)|_L=|x|_K$.若关于 $|\cdot|$ 所诱导距离的每个 Cauchy 列都在 $K$ 中有极限，则称赋值域 $(K,|\cdot|)$ 完备.
\begin{definition*}
赋值域 $(K,|\cdot|)$ 的一个\newterm{nt-completion}{完备化}{completion}是二元组 $((\widehat K,|\cdot|),\iota)$，其中 $(\widehat K,|\cdot|)$ 是完备赋值域，$\iota\colon(K,|\cdot|)\hookrightarrow(\widehat K,|\cdot|)$ 是赋值域态射，满足下述泛性质：对任意完备赋值域 $(L,|\cdot|_L)$ 及任意赋值域态射 $\sigma\colon(K,|\cdot|)\longrightarrow(L,|\cdot|_L)$，存在唯一的 $\widehat\sigma\colon(\widehat K,|\cdot|)\longrightarrow(L,|\cdot|_L)$，使 $\sigma=\widehat\sigma\circ\iota$.
\end{definition*}
换言之，态射 $\iota\colon(K,|\cdot|)\hookrightarrow(\widehat K,|\cdot|)$ 在 $(K,|\cdot|)$ 到完备赋值域的态射中具有泛性质.这个泛性质保证：若完备化存在，则它在唯一同构的意义下唯一.
\begin{proposition*}
每个赋值域 $(K,|\cdot|)$ 都有一个完备化 $((\widehat K,|\cdot|),\iota)$.此外：
\begin{enumerate}[label=\roman*)]
\item $\iota$ 是 $(K,|\cdot|)$ 到完备赋值域的态射中像稠密的唯一一个（相差同构）.
\item $\widehat K$ 的拓扑只依赖于 $K$ 的拓扑.更准确地说，$(K,|\cdot|^a)$ 的完备化是 $((\widehat K,|\cdot|^a),\iota)$.
\end{enumerate}
\end{proposition*}
\begin{proof}
$\widehat K$ 的构造与 $\mathbb R$ 的构造相同.令
\begin{align*}
\widehat K:=\{K\text{ 中的 Cauchy 列}\}/{\sim},\qquad
(x_n)_n\sim(y_n)_n\ \Longleftrightarrow\ \lim_{n\to\infty}|x_n-y_n|=0,
\end{align*}
并以 $|(x_n)_n|:=\lim_{n\to\infty}|x_n|$ 定义范数.态射 $\iota$ 将 $x$ 映到常值列 $(x_n=x)_n$.容易验证所述泛性质成立，且 $K$ 在 $\widehat K$ 中稠密.若 $(K',|\cdot|')$ 是一个完备赋值域，且 $K$ 在其中稠密，则由泛性质得到的态射 $(\widehat K,|\cdot|)\longrightarrow(K',|\cdot|')$ 是同构.最后，ii) 来自两个等价范数具有相同 Cauchy 列这一事实.
\end{proof}
\begin{example*}
$(\mathbb Q,|\cdot|_\infty)$ 的完备化是带自然包含映射的 $(\mathbb R,|\cdot|)$.而 $(\mathbb Q,|\cdot|_p)$ 的完备化记为 $(\mathbb Q_p,|\cdot|_p)$，称为\newterm{nt-p-adic-number-field}{\texorpdfstring{$p$}{p}-进数域}{field of p-adic numbers}.我们将在后面详细研究它.
\end{example*}
\subsection{定理}\label{sec:nt-4-1-5}
\begin{theorem*}
数域 $K$ 上的每个非平凡绝对值，都等价于某个嵌入 $\sigma\colon K\hookrightarrow\mathbb C$ 对应的绝对值 $|\cdot|_\sigma$，或者等价于某个 $\mathfrak p\in\operatorname{Max}(\mathcal O_K)$ 对应的绝对值 $|\cdot|_{\mathfrak p}$.此外，这些绝对值中，只有 $|\cdot|_\sigma$ 与 $|\cdot|_{c\circ\sigma}$ 彼此等价，其中 $c$ 表示复共轭.
\end{theorem*}
\begin{proof}
设 $|\cdot|$ 是 $K$ 上的非平凡绝对值.
\textup{i)} 假设 $|\cdot|$ 是非 Archimedes 的.只需证明，对任意 $x\in\mathcal O_K$，有 $|x|\le1$.事实上，此时令 $\mathfrak p:=\{x\in\mathcal O_K,\ |x|<1\}$，并选取 $\mathcal O_{K,\mathfrak p}$ 的素元 $\varpi$，便有 $|x|=|\varpi|^{v_{\mathfrak p}(x)}$，对所有 $x\in K$ 成立.
设 $x\in\mathcal O_K$，其极小多项式为 $f_x=X^n+a_0X^{n-1}+\cdots+a_0\in\mathbb Z[X]$.由于 $|\cdot|$ 非 Archimedes ，对每个 $i$ 都有 $|a_i|\le1$.于是 $|x|^n\le\max(1,|x|^{n-1})$，从而 $|x|\le1$.
\textup{ii)} 现在假设 $|\cdot|$ 是Archimedes 的.设 $\iota\colon K\hookrightarrow\widehat K$ 为 $K$ 关于 $|\cdot|$ 的完备化.由于 $|\cdot|_{\mathbb Q}\sim|\cdot|_\infty$，包含映射 $\mathbb Q\hookrightarrow K$ 诱导出连续嵌入 $\mathbb R\hookrightarrow\widehat K$.取 $\alpha\in K$ 使 $K=\mathbb Q[\alpha]$.考虑 $\widehat K$ 的子域 $\mathbb R[\iota(\alpha)]$.它在 $\mathbb R$ 上代数，因而次数至多为 $2$.$\widehat K$ 的绝对值在 $\mathbb R[\iota(\alpha)]$ 上诱导一个绝对值 $|\cdot|$，它也是这个 $\mathbb R$ 向量空间上的范数；因此 $\mathbb R[\iota(\alpha)]$ 完备，因为它在 $\mathbb R$ 上有限维，且绝对值 $|\cdot|$ 在 $\mathbb R[\iota(\alpha)]$ 上诱导有限维实向量空间的拓扑，故等价于 $\mathbb R[\iota(\alpha)]$ 上的通常绝对值（该域同构于 $\mathbb R$ 或 $\mathbb C$）.由于完备，$\mathbb R[\iota(\alpha)]$ 在 $\widehat K$ 中闭；而它包含 $\iota(K)$，后者稠密，所以最终得到 $\mathbb R[\iota(\alpha)]=\widehat K$.
我们证明了扩张 $\mathbb R\hookrightarrow\widehat K$ 的次数为 $1$ 或 $2$.若次数为 $1$，则它是拓扑同构；与 $\iota$ 复合，便得到所求的连续嵌入 $\sigma\colon K\hookrightarrow\mathbb R$.若次数为 $2$，则可选取连续的 $\mathbb R$-线性同构 $\widehat K\xrightarrow{\sim}\mathbb C$；与 $\iota$ 复合，便得到所求的连续嵌入 $\sigma\colon K\hookrightarrow\mathbb C$.
\textup{iii)} 还需研究这些范数之间的等价关系.在非 Archimedes 情形，由 $|\cdot|_{\mathfrak p}$ 可以恢复 $\mathfrak p$，因为 $\mathfrak p=\{x\in\mathcal O_K,\ |x|<1\}$.在Archimedes 情形，$\sigma$ 与 $c\circ\sigma$ 显然在 $K$ 上诱导相同的拓扑，因为 $c\colon\mathbb C\longrightarrow\mathbb C$ 连续.反过来，假设 $|\cdot|_\sigma\sim|\cdot|_{\sigma'}$.这意味着，$K$ 关于 $|\cdot|_\sigma$ 的完备化 $K_\sigma$ 与关于 $|\cdot|_{\sigma'}$ 的完备化 $K_{\sigma'}$ 之间，存在一个 $K$-线性的拓扑域同构.注意，$K_\sigma$ 可以等同于 $\sigma(K)$ 在 $\mathbb C$ 中的闭包.于是有图
\begin{align*}
\begin{tikzcd}[ampersand replacement=\&,column sep=large,row sep=large]
K \arrow[r,hook,"\sigma"] \arrow[dr,hook,"\sigma'"']
\& K_\sigma \arrow[r,hook] \arrow[d,dashed,"\exists f\simeq"]
\& \mathbb C \\
\& K_{\sigma'} \arrow[r,hook] \& \mathbb C
\end{tikzcd}
\end{align*}
分为两种情形：
\begin{enumerate}[label=(\alph*)]
\item 若 $K_\sigma=\mathbb R$，则也必有 $K_{\sigma'}=\mathbb R$；但 $\mathbb R$ 唯一的连续自同构是恒等映射，所以 $\sigma=\sigma'$.
\item 若 $K_\sigma=\mathbb C$，则也必有 $K_{\sigma'}=\mathbb C$；但 $\mathbb C$ 的连续自同构只有恒等映射和复共轭，所以 $\sigma=\sigma'$ 或 $\sigma'=c\circ\sigma$.
\end{enumerate}
\end{proof}
\subsection{数域的位}\label{sec:nt-4-1-6}
由前一定理，有
\begin{align*}
\{K\text{ 的位}\}\longleftrightarrow
\operatorname{Max}(\mathcal O_K)\cup\{\text{嵌入 }K\hookrightarrow\mathbb C\}/{\text{共轭}}.
\end{align*}
位通常记为 $v$，相应的 $K$ 的完备化记为 $K_v$.非 Archimedes 位也称为\newterm{nt-finite-place}{有限位}{finite place}.若 $v$ 对应 $\mathfrak p=\mathfrak p_v\mid(p)$，则记作 $v\mid p$.此时也称 $v$ 为一个 $p$-进位，并记 $K_{\mathfrak p}=K_v$.Archimedes 位也称为\newterm{nt-infinite-place}{无限位}{infinite place}，记作 $v\mid\infty$.当 $v$ 对应嵌入 $\sigma$ 时，也记 $K_\sigma=K_v$.若 $K_v\simeq\mathbb R$，则称这个位为实位，否则称为复位.对每个位，仔细选取一个规范化的绝对值十分有用：
\begin{itemize}
\item 若 $v$ 为有限位，取 $|\cdot|_v=|\cdot|_{\mathfrak p_v}$.
\item 若 $v$ 为实位，对应 $\sigma\colon K\hookrightarrow\mathbb R$，取 $|\cdot|_v=|\cdot|_\sigma$.
\item 若 $v$ 为复位，对应 $\{\sigma,c\circ\sigma\}$，取 $|\cdot|_v=|\cdot|_\sigma\times|\cdot|_{c\circ\sigma}$.
\end{itemize}
在最后一种情形，$|\cdot|_v$ 实际上不是绝对值，因为它不满足三角不等式.但这个选择能使我们得到乘积公式.
\subsection{命题}\label{sec:nt-4-1-7}
\begin{proposition*}
采用上述规范化，若 $\alpha\in K$，则有\footnote{译注：最后的乘积公式须要求 $\alpha\ne0$；原文没有写出这一限制.}
\begin{align*}
\prod_{v\mid\infty}|\alpha|_v&=|N_{K/\mathbb Q}(\alpha)|_\infty,
&\prod_{v\mid p}|\alpha|_v&=|N_{K/\mathbb Q}(\alpha)|_p,
&\text{从而}\quad\prod_v|\alpha|_v&=1.
\end{align*}
\end{proposition*}
\begin{proof}
\textup{i)} 由定义，
\begin{align*}
\prod_{v\mid\infty}|\alpha|_v=\prod_{\sigma\colon K\hookrightarrow\mathbb C}|\sigma(\alpha)|=|N_{K/\mathbb Q}(\alpha)|.
\end{align*}
\textup{ii)} 有
\begin{align*}
\prod_{v\mid p}|\alpha|_v
=\prod_{\mathfrak p\mid p}N\mathfrak p^{-v_{\mathfrak p}(\alpha)}
=|\mathcal O_{K,p}/(\alpha)|^{-1}
\end{align*}
（这里 $|\cdot|$ 表示``基数''，这个等式来自分解 $\mathcal O_{K,p}/(\alpha)=\prod_{\mathfrak p\mid p}\mathcal O_{K,p}/\mathfrak p^{v_{\mathfrak p}(\alpha)}$）.
而 $|\mathcal O_{K,p}/(\alpha)|=|(\mathcal O_K/(\alpha))_p|$ 是 $|\mathcal O_K/(\alpha)|=|N_{K/\mathbb Q}(\alpha)|$ 的 $p$-准素部分，因此 $|\mathcal O_{K,p}/(\alpha)|^{-1}=|N_{K/\mathbb Q}(\alpha)|_p$.
\end{proof}
\begin{remark*}
\hyperref[item:nt-2-1-6-real]{原书第14页的同构ii)}可以重新解释为
\begin{align*}
K\otimes_{\mathbb Q}\mathbb R\simeq\prod_{v\mid\infty}K_v.
\end{align*}
有限位也有类似的同构.设 $p$ 为素数.若 $\mathfrak p\mid p$、$\mathfrak p\in\operatorname{Max}(\mathcal O_K)$，则由完备化的泛性质，包含映射 $\mathbb Q\hookrightarrow K$ 诱导出连续嵌入 $\mathbb Q_p\hookrightarrow K_{\mathfrak p}$，因为 $|\cdot|_{\mathfrak p}$ 在 $\mathbb Q$ 上的限制与 $|\cdot|_p$ 等价.因此有典范态射
\begin{align*}
K\otimes_{\mathbb Q}\mathbb Q_p\longrightarrow K_{\mathfrak p},\qquad x\otimes y\longmapsto xy.
\end{align*}
若写成 $K=\mathbb Q[\alpha]$，便看到 $K_{\mathfrak p}$ 的范数在有限维 $\mathbb Q_p$ 向量空间 $\mathbb Q_p[\iota(\alpha)]\subset K_{\mathfrak p}$ 上诱导一个范数，因而后者完备.由于它包含 $K$，有 $\mathbb Q_p[\iota(\alpha)]=K_{\mathfrak p}$，所以 $K_{\mathfrak p}$ 在 $\mathbb Q_p$ 上有限维.现在考虑态射
\begin{align*}
K\otimes_{\mathbb Q}\mathbb Q_p\longrightarrow\prod_{\mathfrak p\mid p}K_{\mathfrak p}.
\end{align*}
逼近引理（定理~\ref{sec:nt-3-3-4}的推论 iii)）说明，$K$ 在右端的像稠密（习题）.因此，上述态射的像是右端的一个稠密 $\mathbb Q_p$ 向量子空间，故等于整个右端，因为完备赋值域上的每个有限维向量空间都完备.于是，\emph{上述乘积态射是满射}.
我们将在后面证明这个态射是同构.事实上，$K_{\mathfrak p}$ 的显式描述将说明两端在 $\mathbb Q_p$ 上的维数相同.
\end{remark*}
\subsection{相对情形}\label{sec:nt-4-1-8}
设给定数域扩张 $L\supset K$.通过限制，可得到从 $L$ 的位集合到 $K$ 的位集合的映射 $w\mapsto v$.记作 $w\mid v$，读作``$w$ 整除 $v$''或``$w$ 位于 $v$ 之上''.当然，此时 $w$ 为非 Archimedes 位当且仅当 $v$ 为非 Archimedes 位，并且有 $\mathfrak P_w\mid\mathfrak p_v$.前一命题的公式可推广为（见后文）：
\begin{align*}
\forall\alpha\in L,\qquad\prod_{w\mid v}|\alpha|_w&=|N_{L/K}(\alpha)|_v,\\
L\otimes_K K_v&\xrightarrow{\sim}\prod_wL_w.
\end{align*}

\section{非 Archimedes 情形的完备化}\label{sec:nt-4-2}
\subsection{一般性质}\label{sec:nt-4-2-1}
设 $(K,|\cdot|)$ 是非 Archimedes 赋值域，$\mathcal O$ 为它的赋值环，$\mathfrak m$ 为该环的极大理想.
它的完备化 $((\widehat K,|\cdot|),\iota)$ 也是非 Archimedes 的（例如，因为 $|\cdot|$ 在整数上仍有界）.记它的赋值环为 $\widehat{\mathcal O}$，极大理想为 $\widehat{\mathfrak m}$.
于是 $\iota(\mathcal O)\subset\widehat{\mathcal O}$，且 $(\widehat{\mathcal O},\iota)$ 是 $\mathcal O$ 关于绝对值 $|\cdot|$ 所诱导的度量的一个完备化.
现在假设 $|K^\times|$ 是离散的（等价地，$\mathcal O$ 是离散赋值环）.这时有以下性质：
\begin{enumerate}[label=\roman*)]
\item $|\widehat K^\times|=|K^\times|$（由 Cauchy 列的构造可直接看出），所以 $\widehat{\mathcal O}$ 也是离散赋值环，且 $\mathfrak m$ 的任一素元 $\varpi$ 也是 $\widehat{\mathfrak m}$ 的素元（因为 $\varpi$ 是素元当且仅当 $|\varpi|$ 生成 $|K^\times|$）.
\item $\mathcal O$ 上与 $|\cdot|$ 相伴的拓扑是由形如 $x+\mathfrak m^i=x+\varpi^i\mathcal O$ 的集合生成的``$\varpi$ 进拓扑''，其中 $x\in\mathcal O$、$i\in\mathbb N$.事实上，对任意 $1>\varepsilon>0$，$x+\mathfrak m^i$ 就是开球 $B(x,|\varpi|^{i+\varepsilon})$.\footnote{译注：原文球半径指数为 $i+\varepsilon$，照录.由于 $0<|\varpi|<1$，这一半径小于 $|\varpi|^i$，与所述球等式有冲突.}
\item 同样，$\widehat{\mathcal O}$ 上与 $|\cdot|$ 相伴的拓扑是由形如 $x+\widehat{\mathfrak m}^{\,i}=x+\varpi^i\widehat{\mathcal O}$ 的集合生成的``$\varpi$ 进拓扑''，其中 $x\in\mathcal O$、$i\in\mathbb N$.
\item 包含映射 $\iota$ 诱导同构 $\mathcal O/\varpi^n\mathcal O\xrightarrow{\sim}\widehat{\mathcal O}/\varpi^n\widehat{\mathcal O}$.事实上，
\begin{itemize}
\item 有 $\varpi^n\widehat{\mathcal O}\cap\mathcal O=\{x\in\mathcal O,\ |x|\le|\varpi^n|\}=\varpi^n\mathcal O$，由此得到单射性.
\item $\mathcal O$ 在 $\widehat{\mathcal O}$ 中关于 $\varpi$ 进拓扑稠密，由此得到满射性.
\end{itemize}
\end{enumerate}
现在考虑环
\begin{align*}
\varprojlim_n\mathcal O/\varpi^n\mathcal O
:=\left\{(x_n)_{n\in\mathbb N}\in\prod_n\mathcal O/\varpi^n\mathcal O,\quad x_n=\overline{x_m},\ \forall m>n\right\}.
\end{align*}
在每个 $\mathcal O/\varpi^n\mathcal O$ 上取离散拓扑，再赋予此环由乘积拓扑诱导的拓扑.这样得到一个拓扑环（加法和乘法连续）.以上第 iv) 项给出环同态 $\widehat{\mathcal O}\longrightarrow\mathcal O/\varpi^n\mathcal O$；由第 iii) 项，它们对 $\mathcal O/\varpi^n\mathcal O$ 的离散拓扑连续.这些同态的笛卡儿积于是诱导一个连续的 $\mathcal O$ 线性拓扑环同态
\begin{align*}
\widehat{\mathcal O}\longrightarrow\varprojlim_n\mathcal O/\varpi^n\mathcal O,
\qquad x\longmapsto(\overline x)_n.
\end{align*}
\subsection{命题}\label{sec:nt-4-2-2}
\begin{proposition*}
在上述假设下（$|K^\times|$ 在 $\mathbb R_{>0}^\times$ 中离散），以上同态是拓扑环同构
\begin{align*}
\widehat{\mathcal O}\xrightarrow{\sim}\varprojlim_n\mathcal O/\varpi^n\mathcal O.
\end{align*}
\end{proposition*}
\begin{proof}
\emph{单射性.}由以上第 iv) 项，可以将所研究同态的目标换成 $\varprojlim_n\widehat{\mathcal O}/\varpi^n\widehat{\mathcal O}$.于是此同态的核为 $\bigcap_n\varpi^n\widehat{\mathcal O}$；因为 $\varpi$ 是离散赋值环 $\widehat{\mathcal O}$ 的素元，这个核为零.
\emph{满射性.}设 $(x_n)_n$ 是右端的一个序列，$(\widetilde x_n)_n$ 是它在 $\mathcal O^{\mathbb N}$ 中的一列提升，则 $(\widetilde x_n)$ 是 $\mathcal O$ 中关于 $|\cdot|$ 所对应度量的 Cauchy 列.因此它在 $\widehat{\mathcal O}$ 中有极限 $x$，而 $x$ 的像恰好为 $(x_n)_n$.
\emph{同胚性.}连续性已作说明.只须证明这个同态是开映射.它将 $x+\mathfrak m^i$ 映到投影
\begin{align*}
\varprojlim_n(\widehat{\mathcal O}/\varpi^n\widehat{\mathcal O})\longrightarrow\widehat{\mathcal O}/\varpi^i\widehat{\mathcal O}
\end{align*}
下 $\{\overline x\}$ 的原像.这个投影连续，而 $\{\overline x\}$ 在 $\widehat{\mathcal O}/\varpi^i\widehat{\mathcal O}$ 中是开集，所以 $x+\mathfrak m^i$ 的像确实是开集.\footnote{译注：本段的 $x+\mathfrak m^i$ 沿用原文；同态的定义域为 $\widehat{\mathcal O}$，通常应在这里使用 $x+\widehat{\mathfrak m}^{\,i}$.}
\end{proof}
这个命题给出了完备化的一种代数构造.具体地说，若 $R$ 是自然投影的一个集合截面\linebreak
$\mathcal O/\varpi\mathcal O\hookrightarrow\mathcal O$ 的像，则 $\mathcal O$ 中每个元素都唯一地写成级数
\begin{align*}
x=\sum_{n\in\mathbb N}a_n\varpi^n.
\end{align*}
其中 $a_n\in R$.\footnote{译注：原文这里写的是 $\mathcal O$，照录；上述所有级数所描述的通常是 $\widehat{\mathcal O}$ 中的元素.}不过，用这种表示不易描述加法和乘法.
\begin{example*}
记 $\mathbb Q_p$ 的赋值环为 $\mathbb Z_p$.它就是 $\mathbb Z$ 的``$p$ 进完备化''，由下式给出：
\begin{align*}
\mathbb Z_p=\varprojlim_n\mathbb Z/p^n\mathbb Z
=\left\{\sum_{n\in\mathbb N}a_np^n,\ (a_n)_n\in\{0,\ldots,p-1\}^{\mathbb N}\right\}.
\end{align*}
注意，$\mathbb N$ 可视为那些在 $n\gg0$ 时满足 $a_n=0$ 的元素组成的集合，它在 $\mathbb Z_p$ 中稠密.通常的 $p$ 进位制乘法和加法规则（带进位等）延拓到 $\mathbb Z_p$.例如，$-1$ 写成 $-1=\sum_{n\in\mathbb N}(p-1)p^n$.此外，因为 $\mathbb Q_p=\mathbb Z_p[\frac1p]$，也得到如下``级数''展开：
\begin{align*}
\mathbb Q_p=\left\{\sum_{n\in\mathbb Z}a_np^n,\ (a_n)_n\in\{0,\ldots,p-1\}^{\mathbb N}\text{ 且 }a_{-n}=0\text{ 对 }n\gg0\right\}.
\end{align*}
\end{example*}
\begin{example*}
更容易的例子是 $\mathbb F_p[X]$ 上的 $X$ 进赋值.这时完备化为
\begin{align*}
\mathbb F_p[[X]]=\varprojlim_n\mathbb F_p[X]/(X^n)
=\left\{\sum_{n\in\mathbb N}a_nX^n,\ (a_n)_n\in\mathbb F_p^{\mathbb N}\right\}.
\end{align*}
不过，这里的加法和乘法容易理解得多：它们就是形式幂级数上通常的运算.分式域记为 $\mathbb F_p((X))$.
\end{example*}
\subsection{推论}\label{sec:nt-4-2-3}
\begin{corollary*}
若 $K$ 是非 Archimedes 完备赋值域，则以下条件等价：
\begin{enumerate}[label=\roman*)]
\item $K$ 局部紧.
\item $\mathcal O$ 紧.
\item $|K^\times|$ 在 $\mathbb R^\times$ 中离散，且剩余域 $\mathcal O/\mathfrak m$ 有限.
\end{enumerate}
\end{corollary*}
\begin{proof}
$i)\Rightarrow ii)$：若 $K$ 局部紧，它的闭球都紧，而 $\mathcal O$ 是一个闭球.
$iii)\Rightarrow i)$ 由上一命题得到，因为所有 $\mathcal O/\varpi^n\mathcal O$ 都有限，所以它们的乘积紧，而投射极限在乘积中是闭的.
$ii)\Rightarrow iii)$：由于 $\mathfrak m$ 是开集，投影 $\mathcal O\longrightarrow\mathcal O/\mathfrak m$ 对 $\mathcal O/\mathfrak m$ 的离散拓扑连续.它又是满射，所以 $\mathcal O/\mathfrak m$ 既紧又离散，因而有限.现在取 $0<r<1$ 并令 $I_r:=\{x\in\mathcal O,\ |x|<r\}$.这是一个开理想，因此如上所述，$\mathcal O/I_r$ 有限.这蕴含 $|K^\times|\cap\mathopen]r,1\mathclose]$ 有限，从而 $\mathbb R_{>0}^\times$ 的子群 $|K^\times|$ 离散.
\end{proof}
\begin{example*}
环 $\mathbb Z_p$ 紧，$\mathbb Q_p$ 局部紧.更一般地，若 $K$ 为数域且 $\mathfrak p\mid p$，则 $K_{\mathfrak p}$ 是局部紧域，且为 $\mathbb Q_p$ 的有限次扩张（上文已经见过）.同样，$\mathbb F_p((X))$ 局部紧.
\end{example*}
这些非 Archimedes 局部紧域的例子实际上具有代表性：
\begin{proposition*}
每个非 Archimedes 局部紧域都是某个 $\mathbb Q_p$ 的有限扩张，或者同构于某个 $\mathbb F_q[[X]]$.\footnote{译注：原文为 $\mathbb F_q[[X]]$，照录；这一幂级数环并非域，通常的分类结论使用 Laurent 级数域 $\mathbb F_q((X))$.}
\end{proposition*}
我们稍后会证明，反过来，$\mathbb Q_p$ 的每个有限扩张都有一个使其成为局部紧赋值域的绝对值；也会证明这样的有限扩张来自某个数域的完备化.
\begin{proof}
设 $K$ 的特征为 $0$，$p$ 为有限域 $\mathcal O/\mathfrak m$ 的特征.则 $|p|<1$，且 $|\cdot|$ 在 $\mathbb Q$ 上的限制是 $p$ 进绝对值，因而 $K$ 包含 $\mathbb Q_p$，$\mathcal O$ 包含 $\mathbb Z_p$.我们要证明 $\mathcal O$ 在 $\mathbb Z_p$ 上有限生成.已知 $\mathcal O/p\mathcal O$ 有限.选定 $\mathcal O/p\mathcal O$ 的一个 $\mathbb F_p$ 基 $\overline x_1,\ldots,\overline x_m$，并在 $\mathcal O$ 中取提升 $x_1,\ldots,x_m$.我们要证明 $x_1,\ldots,x_m$ 是 $\mathcal O$ 的一个 $\mathbb Z_p$ 基.
若存在非零的 $\mathbb Z_p$ 线性相关关系 $\sum_i a_ix_i=0$，则存在一个这样的关系使 $v_p(a_1)=0$.但投影到 $\mathcal O/p\mathcal O$ 就得到 $\overline x_i$ 之间的一个非零线性相关关系，矛盾.因此族 $(x_1,\ldots,x_n)$ 在 $\mathbb Z_p$ 上线性无关.
现在取 $x\in\mathcal O$.可找到 $a_1^1,\ldots,a_m^1\in\mathbb Z_p$ 使 $x-\sum_i a_i^1x_i\in p\mathcal O$.递归地，可找到 $a_1^n,\ldots,a_m^n\in\mathbb Z_p$，使 $x-\sum_i a_i^nx_i\in p^n\mathcal O$ 且 $v_p(a_i^n-a_i^{n-1})\ge n$.因此，序列 $(a_i^n)_{n\in\mathbb N}$ 都是 Cauchy 列，收敛到某个 $a_i\in\mathbb Z_p$.由构造，
\begin{align*}
x-\sum_i a_ix_i\in\bigcap_n p^n\mathcal O
=\{x\in\mathcal O_K,\ |x|<|p|^n\ \forall n\}=\{0\}.
\end{align*}
所以族 $(x_1,\ldots,x_m)$ 是 $\mathcal O$ 在 $\mathbb Z_p$ 上的一个基，因而也是 $K$ 在 $\mathbb Q_p$ 上的一个基（习题）.
$K$ 为正特征的情形留给读者.
\end{proof}
\subsection{Dedekind 环的完备化与扩张}\label{sec:nt-4-2-4}
回到现在已经熟悉的情形：设 $A$ 为 Dedekind 环，$K=\operatorname{Frac}(A)$，$L$ 为 $K$ 的有限可分扩张，$B$ 为 $A$ 在 $L$ 中的整闭包.我们知道 $B$ 作为 $A$ 模有限生成.最后，取 $\mathfrak p\in\operatorname{Max}(A)$.我们已见到，$B_{\mathfrak p}:=A_{\mathfrak p}\otimes_A B$ 通过映射 $a\otimes b\longmapsto ab\in L$ 与 $L$ 的子环 $A_{\mathfrak p}B$ 等同，而后者正是 $A_{\mathfrak p}$ 在 $L$ 中的整闭包.这是一个半局部环，且有
\begin{align*}
\operatorname{Max}(B_{\mathfrak p})\longleftrightarrow
\{\mathfrak P\in\operatorname{Max}(B),\ \mathfrak P\mid\mathfrak p\}.
\end{align*}
这个环是整环，但下面的结果表明，其 $\mathfrak p$ 进完备化可以将它分裂为若干完备局部环的乘积.
\begin{theorem*}
沿用上述记号，有拓扑环同构
\begin{align*}
\widehat A_{\mathfrak p}\otimes_A B\xrightarrow{\sim}
\prod_{\mathfrak P\mid\mathfrak p}\widehat B_{\mathfrak P}.
\end{align*}
此外，$\widehat B_{\mathfrak P}$ 是秩为 $e(\mathfrak P/\mathfrak p)f(\mathfrak P/\mathfrak p)$ 的自由 $\widehat A_{\mathfrak p}$ 模，并可视为 $\widehat A_{\mathfrak p}$ 在 $L$ 关于任一 $\mathfrak P$ 进绝对值的完备化 $\widehat L_{\mathfrak P}$ 中的整闭包.
\end{theorem*}
\begin{proof}
按定义，左端为
\begin{align*}
\widehat A_{\mathfrak p}\otimes_A B
&=\left(\varprojlim_n A/\mathfrak p^n\right)\otimes_A B\\
&=\left(\varprojlim_n A_{\mathfrak p}/\mathfrak p^nA_{\mathfrak p}\right)\otimes_{A_{\mathfrak p}} B_{\mathfrak p}.
\end{align*}
因为 $B_{\mathfrak p}$ 是有限秩自由模，可以将它移入投射极限中（因为投射系统直和的投射极限是各投射极限的直和）.于是左端变为
\begin{align*}
\varprojlim_n(A_{\mathfrak p}/\mathfrak p^nA_{\mathfrak p}\otimes_{A_{\mathfrak p}} B_{\mathfrak p})
=\varprojlim_n(A/\mathfrak p^n\otimes_A B)
=\varprojlim_n B/\mathfrak p^nB.
\end{align*}
但对每个 $n$，用通常的记号，有
\begin{align*}
B/\mathfrak p^nB=\prod_{\mathfrak P\mid\mathfrak p}B/\mathfrak P^{ne(\mathfrak P/\mathfrak p)}.
\end{align*}
取投射极限就得到定理中的同构，这里用到了以下事实：若 $[(M_n)_n,(\rho_{m,n}\colon M_m\longrightarrow M_n)_{n\le m}]$ 是一个投射系统，$\varphi\colon\mathbb N\longrightarrow\mathbb N$ 是严格递增映射，则
\begin{align*}
\varprojlim_n M_n\xrightarrow{\sim}\varprojlim_n M_{\varphi(n)}.
\end{align*}
$\widehat B_{\mathfrak P}$ 在 $\widehat A_{\mathfrak p}$ 上自由，因为它是自由模的直和因子，而 $\widehat A_{\mathfrak p}$ 是主理想环（离散赋值环）.其秩就是
\begin{align*}
\widehat B_{\mathfrak P}/\mathfrak p\widehat B_{\mathfrak P}
=\widehat B_{\mathfrak p}/\mathfrak P^{e(\mathfrak P/\mathfrak p)}\widehat B_{\mathfrak P}
=B_{\mathfrak p}/\mathfrak P^{e(\mathfrak P/\mathfrak p)}B_{\mathfrak P}
=B/\mathfrak P^{e(\mathfrak P/\mathfrak p)}
\end{align*}
在 $A/\mathfrak p$ 上的秩，即 $e(\mathfrak P/\mathfrak p)f(\mathfrak P/\mathfrak p)$.\footnote{译注：上式中的 $\widehat B_{\mathfrak p}$、$B_{\mathfrak p}$ 沿用原文下标，原文在这条等式中混用了 $\mathfrak p$ 和 $\mathfrak P$.}
现在赋予 $L$ 一个 $\mathfrak P$ 进绝对值（形如 $|x|=\rho^{-v_{\mathfrak P}(x)}$，其中 $\rho>1$）.它在 $K$ 上的限制是 $\mathfrak p$ 进绝对值，因而有完备化之间的连续嵌入 $\widehat K_{\mathfrak p}\hookrightarrow\widehat L_{\mathfrak P}$.已经见过，$\widehat A_{\mathfrak p}$ 可视为 $\widehat K_{\mathfrak p}$ 的赋值环，$\widehat B_{\mathfrak P}$ 在 $\widehat L_{\mathfrak P}$ 中也同样如此.特别地，$\widehat B_{\mathfrak P}$ 是离散赋值环，因而正规，且其分式域为 $\widehat L_{\mathfrak P}$.因为它在 $\widehat A_{\mathfrak p}$ 上有限生成，所以是后者在 $\widehat L_{\mathfrak P}$ 中的整闭包.
\end{proof}
\begin{example*}
设 $K$ 为数域，则有 $\mathcal O_K\otimes\mathbb Z_p\xrightarrow{\sim}\prod_{\mathfrak p\mid p}\mathcal O_{K_{\mathfrak p}}$.将 $p$ 逆化，得到同构
\begin{align*}
K\otimes_{\mathbb Q}\mathbb Q_p\xrightarrow{\sim}\prod_{\mathfrak p\mid p}K_{\mathfrak p}.
\end{align*}
更一般地，若 $L\supset K$ 是数域扩张，$v$ 是 $K$ 的一个有限位，则得到分解
\begin{align*}
L\otimes_K K_v\xrightarrow{\sim}\prod_{w\mid v}L_w.
\end{align*}
\end{example*}
\subsection{赋值的延拓}\label{sec:nt-4-2-5}
我们考察以下问题：给定赋值域 $(K,|\cdot|)$ 及 $K$ 的扩张 $L$，能否将 $K$ 的绝对值延拓到 $L$？我们在 $K$ 非 Archimedes、完备且 $|K^\times|$ 离散的假设下回答这个问题.结果既可以用代数方式，也可以用更``分析''的方式陈述.下面先给出代数表述.
\begin{corollary*}
设 $A,K,B,L$ 如上一分目.若 $A$ 是完备离散赋值环，则 $B$ 也是完备离散赋值环.
\end{corollary*}
\begin{proof}
上一定理中的同构成为
\begin{align*}
B\xrightarrow{\sim}\prod_{\mathfrak P\in\operatorname{Max}(B)}\widehat B_{\mathfrak P}.
\end{align*}
但 $B$ 是整环，故 $\operatorname{Max}(B)$ 是单点集 $\{\mathfrak P\}$，且 $B=\widehat B_{\mathfrak P}$ 确为完备离散赋值环.
\end{proof}
下面是更偏分析的版本.
\begin{corollary*}
设 $(K,|\cdot|_K)$ 是非 Archimedes 完备赋值域，$|K^\times|$ 在 $\mathbb R_{>0}^\times$ 中离散，且 $L$ 为一个次数为 $n$ 的有限可分扩张.则在 $L$ 上存在唯一一个延拓 $|\cdot|_K$ 的绝对值 $|\cdot|_L$，由下式给出：
\begin{align*}
\forall x\in L,\qquad |x|_L=|N_{L/K}(x)|_K^{1/n}.
\end{align*}
\end{corollary*}
\begin{proof}
将上一推论应用于 $A=\mathcal O=\{x\in K,\ |x|\le1\}$，得到延拓的存在性.事实上，写 $|x|_K=\rho^{-v(x)}$，其中 $v$ 是 $A$ 的赋值，$\rho=|\varpi|^{-1}$，$\varpi$ 为 $A$ 的一个素元.设 $v'$ 为 $A$ 在 $L$ 中的整闭包 $B$ 的赋值.则 $x\longmapsto\rho^{-v'(x)/v'(\varpi)}$ 定义了 $L$ 上延拓 $|\cdot|_K$ 的绝对值.
下面证明唯一性.设 $|\cdot|_L$ 是一个延拓.它是一个非 Archimedes 绝对值，其赋值环包含 $A$，因而也包含 $A$ 的整闭包 $B$.所以 $|B^\times|_L=1$.由分解 $L^\times=\bigsqcup_n\varpi_B^nB^\times$ 可见，$|\cdot|_L$ 由 $|\varpi_B|_L$ 确定.取 $e$ 使 $\varpi B=\varpi_B^eB$（即 $e=v'(\varpi)$），则必须有 $|\varpi_B|_L^e=|\varpi|_K$，这就完全确定了 $|\cdot|_L$.
为证明所述公式，可以直接证明 $v(N_{L/K}(x))=\frac ne v'(x)$（习题）.也可以选定 $L$ 在 $K$ 上的一个 Galois 闭包 $\widetilde L$，并注意到，由唯一性，对每个嵌入 $\sigma\colon L\hookrightarrow\widetilde L$，都有 $|x|_L=|\sigma(x)|_{\widetilde L}$，从而
\begin{align*}
|x|_L^n=\prod_{\sigma\colon L\hookrightarrow\widetilde L}|\sigma(x)|_{\widetilde L}
=|N_{L/K}(x)|_K.
\end{align*}
\end{proof}
\emph{应用.}设 $\overline{\mathbb Q}_p$ 为 $\mathbb Q_p$ 的一个代数闭包.$|\cdot|_p$ 有唯一的延拓 $|\cdot|_p$ 到 $\overline{\mathbb Q}_p$，且 Galois 群的作用是等距的：
\begin{align*}
\forall\sigma\in G_{\overline{\mathbb Q}_p/\mathbb Q_p},\quad
\forall x\in\overline{\mathbb Q}_p,\qquad |\sigma(x)|_p=|x|_p.
\end{align*}
若 $K$ 为数域，每个嵌入 $\sigma\colon K\hookrightarrow\overline{\mathbb Q}_p$ 都给出 $K$ 的一个 $p$ 进绝对值 $|x|_\sigma:=|\sigma(x)|_p$.两个在 $G_{\overline{\mathbb Q}_p/\mathbb Q_p}$ 作用下共轭的嵌入给出相同的绝对值，这诱导一个双射
\begin{align*}
\{\sigma\colon K\hookrightarrow\overline{\mathbb Q}_p\}/G_{\overline{\mathbb Q}_p/\mathbb Q_p}
\longleftrightarrow\{p\text{ 进位 }v\}.
\end{align*}
完备化 $K_\sigma=K_v$ 通过 $\sigma$ 与 $\overline{\mathbb Q}_p$ 中的复合子域 $\sigma(K)\mathbb Q_p$ 等同，后者也是 $\sigma(K)$ 在 $\overline{\mathbb Q}_p$ 中的闭包.现在的情形与 Archimedes 情形十分相似.
\begin{remark*}
$\overline{\mathbb Q}_p$ 的绝对值不再离散！事实上，$|\overline{\mathbb Q}_p^{\,\times}|_p=p^{\mathbb Q}$.包含关系 $|\overline{\mathbb Q}_p^{\,\times}|_p\subset p^{\mathbb Q}$ 由推论中的公式（含范数的那个公式）得到；另一个包含关系则由以下观察得到：若 $\alpha^n=p$，则 $|\alpha|=p^{-1/n}$.特别地，$\overline{\mathbb Q}_p$ 不完备.它的完备化通常记为 $\mathbb C_p$.可以证明它仍然是代数闭的.
\end{remark*}
\begin{exercise*}
设 $F\supset\mathbb Q_p$ 为有限扩张.若唯一的 $\mathfrak P\mid(p)$ 分歧或不分歧，则分别称此扩张``分歧''或``不分歧''.也记 $e=e(\mathfrak P/p)$（分歧指数）、$f=f(\mathfrak P/p)$（剩余次数）.有 $[F:\mathbb Q_p]=ef$.
\begin{enumerate}[label=\roman*)]
\item 证明 $F$ 不分歧当且仅当 $|F^\times|_p=|\mathbb Q_p^\times|_p$.
\item 证明 $F$ 完全分歧（即 $f=1$）当且仅当对 $\mathcal O_F$ 的每个素元 $\varpi_F$，范数 $N_{F/\mathbb Q_p}(\varpi_F)$ 都是 $\mathbb Z_p$ 的素元.
\end{enumerate}
\end{exercise*}
\emph{注意.}在 $F\supset\mathbb Q_p$ 上，不要将上面得到的绝对值 $|\cdot|_p$ 与按数域情形规范化的绝对值 $|\cdot|_{\mathfrak p_F}$ 混淆.若 $v$ 为 $\mathcal O_F$ 的规范化赋值，则有\footnote{译注：原文先写 $\mathfrak p_F$，公式中则写 $\mathfrak p_E$，照录.}
\begin{align*}
\forall x\in F,\qquad |x|_p=p^{-v(x)/e}
\qquad\text{且}\qquad |x|_{\mathfrak p_E}=p^{-fv(x)}.
\end{align*}
\subsection{Galois 作用}\label{sec:nt-4-2-6}
沿用前面各分目中的 $A,K,B,L$，并进一步假设 $L$ 在 $K$ 上是 Galois 扩张.取 $\mathfrak P\in\operatorname{Max}(B)$，令 $\mathfrak p=A\cap\mathfrak P$.分解群 $D_{\mathfrak P}\subset G_{L/K}$ 作用在每个 $B/\mathfrak P^i$ 上，所以作用在投射极限 $\widehat B_{\mathfrak P}$ 上，进而作用在其分式域 $\widehat L_{\mathfrak P}=\operatorname{Frac}(\widehat B_{\mathfrak P})$ 上；后者也是 $L$ 关于任一 $\mathfrak P$ 进绝对值的完备化.这个作用由域自同构给出，且保持 $A$，因而保持 $\widehat A_{\mathfrak p}$，从而保持 $\widehat K_{\mathfrak p}\subset\widehat L_{\mathfrak P}$.
\begin{proposition*}
由此得到的同态 $D_{\mathfrak P}\longrightarrow G_{\widehat L_{\mathfrak P}/\widehat K_{\mathfrak p}}$ 是双射.
\end{proposition*}
\begin{proof}
因为 $K\subset\widehat K_{\mathfrak p}$、$L\subset\widehat L_{\mathfrak P}$，通过限制也有同态
\begin{align*}
G_{\widehat L_{\mathfrak P}/\widehat K_{\mathfrak p}}
\xrightarrow{\operatorname{res}}G_{L/K}.
\end{align*}
显然，$\operatorname{res}$ 与命题中的同态复合是 $D_{\mathfrak P}$ 上的恒等映射.这证明了命题中同态的单射性.由于 $|D_{\mathfrak P}|=ef=[\widehat L_{\mathfrak P}:\widehat K_{\mathfrak g}]$，便得到满射性.也可看出 $\operatorname{res}$ 是逆同构.\footnote{译注：原文扩张次数分母的下标为 $\mathfrak g$，照录；本段其余地方均为 $\mathfrak p$.}
\end{proof}
\section{Hensel 引理}\label{sec:nt-4-3}
\subsection{定理}\label{sec:nt-4-3-1}
\begin{theorem*}[\newterm{nt-hensel-lemma}{Hensel 引理}{Hensel's lemma}]
设 $A$ 是剩余域为 $k$ 的完备离散赋值环，$f\in A[X]$ 为首一多项式，其在 $k[X]$ 中的像为 $\overline f$.假设在 $k[X]$ 中有分解 $\overline f=g_0h_0$，且 $(g_0,h_0)=1$.则存在唯一一对 $(g,h)\in A[X]^2$，使 $f=gh$、$\overline g=g_0$、$\overline h=h_0$，且 $\deg(g)=\deg(g_0)$、$\deg(h)=\deg(h_0)$.\footnote{译注：原文未要求因子 $g,h$ 首一，照录；若不作归一化，可取约化为 $1$ 的单位 $u$，将两因子分别乘以 $u$ 与 $u^{-1}$，所述唯一性通常不成立.}
\end{theorem*}
\emph{应用举例.}考虑 $A=\mathbb Z_p$ 及 $f(X)=X^{p-1}-1$.剩余多项式 $\overline f$ 在 $\mathbb F_p[X]$ 中完全分裂（且可分）.所以 Hensel 引理蕴含 $f$ 也完全分裂，即 $\mathbb Q_p$ 包含所有 $p-1$ 次单位根.
\begin{remark*}
注意，假设 $(g_0,h_0)=1$ 对提升分解至关重要.例如在 $\mathbb Z_p[X]$ 中，多项式 $f(X)=X^n-p$ 不可约（Eisenstein 判别法），尽管其约化 $\overline f(X)=X^n$ 可约.
\end{remark*}
\subsection{代数证明}\label{sec:nt-4-3-2}
通常的证明是递归构造多项式对 $(g_n,h_n)\in A[X]$，两者的次数分别为 $\deg(g_0)$ 和 $\deg(h_0)$，并满足
\begin{align*}
f-g_nh_n&\in\varpi^nA[X],\\
g_n-g_{n-1}&\in\varpi^nA[X]\quad\text{且}\quad h_n-h_{n-1}\in\varpi^nA[X].
\end{align*}
因为 $A$ 完备，取极限就得到 $g$ 和 $h$ 的存在性.细节可参看 Milne 的讲义，或者作为习题完成.
下面给出一种略有不同的方法.更一般地，设 $A$ 是极大理想为 $\mathfrak m$ 的 Noether 局部环，$B$ 是作为 $A$ 模有限生成的 $A$ 代数.投影 $B\twoheadrightarrow B/\mathfrak m$ 诱导映射
\begin{align*}
\{B\text{ 的幂等元}\}\longrightarrow\{B/\mathfrak mB\text{ 的幂等元}\}.
\end{align*}
若对每个有限 $A$ 代数 $B$，上述映射都是满射，则称 $A$ 为\newterm{nt-henselian-ring}{Hensel 环}{Henselian ring}.容易验证（习题），此时该映射也是双射.[若 $\varepsilon\equiv\varepsilon'\pmod{\mathfrak m}$，则 Nakayama 引理表明 $\varepsilon(1-\varepsilon')=0=\varepsilon'(1-\varepsilon)$.]
\begin{proposition*}
在正规 Hensel 局部环中，Hensel 引理成立.
\end{proposition*}
\begin{proof}
设 $f,g_0,h_0$ 如上面 Hensel 引理的陈述.令 $B:=A[X]/(f)$.于是有 $B/\mathfrak mB=k[X]/(g_0)\times k[X]/(h_0)$（中国剩余定理），从而有幂等元 $\varepsilon$，使 $\varepsilon(B/\mathfrak mB)=k[X]/(g_0)$、$(1-\varepsilon)(B/\mathfrak m)=k[X]/(h_0)$.具体而言，若 $u,v\in k[X]$ 满足 $uh_0+vg_0=1$，则 $\varepsilon$ 是 $uh_0$ 在 $B/\mathfrak mB=k[X]/(\overline f)$ 中的像.反过来，有 $h_0=\gcd(\overline f,\overline f_{\varepsilon})$，其中 $\overline f_{\varepsilon}\in k[X]$ 代表 $\varepsilon\in k[X]/(\overline f)$.
由于 $A$ 为 Hensel 环，可以将 $\varepsilon$ 提升为 $B$ 的一个幂等元 $\widetilde\varepsilon$.设 $f_{\widetilde\varepsilon}\in A[X]$ 在 $B$ 中代表 $\varepsilon$.\footnote{译注：此句原文写 $\varepsilon$，而其上下文讨论的是提升后的 $\widetilde\varepsilon$，照录.}令 $h:=\gcd(f,f_{\widetilde\varepsilon})$ 及 $g:=\gcd(f,1-f_{\widetilde\varepsilon})$（最大公因式在主理想环 $K[X]$ 中取首一的，其中 $K=\operatorname{Frac}(A)$）.显然 $(g,h)=1$ 且 $gh\mid f$；由 $f_{\widetilde\varepsilon}(1-f_{\widetilde\varepsilon})\in(f)$，最终有 $gh=f$.因为 $A$ 正规，$g$ 和 $h$ 都属于 $A[X]$；由构造，$\overline g\mid g_0$ 且 $\overline h\mid h_0$.\footnote{译注：这两个整除方向沿用原文.}但因为 $\overline{gh}=\overline f$，确有 $\overline g=g_0$、$\overline h=h_0$.
对 $(g,h)$ 的唯一性由提升 $\widetilde\varepsilon$ 的唯一性得到.
\end{proof}
\begin{theorem*}
完备 Noether 局部环是 Hensel 环.
\end{theorem*}
这里的``完备''是指关于由 $x+\mathfrak m^i$（$x\in A$、$i\in\mathbb N$）生成的``$\mathfrak m$ 进拓扑''分离且完备.
\begin{proof}
写成 $1=(X+(1-X))^{2n}=\sum_{k=0}^{2n}\binom{2n}{k}X^{2n-k}(1-X)^k$，并令
\begin{align*}
f_n(X):=\sum_{k=0}^n\binom{2n}{k}X^{2n-k}(1-X)^k.
\end{align*}
通过计算可以验证以下性质：
\begin{enumerate}[label=\roman*)]
\item $f_n(X)\equiv f_n(X)^2\pmod{X^n(1-X)^n}$.
\item $f_n(X)\equiv f_{n-1}(X)\pmod{X^{n-1}(1-X)^{n-1}}$.
\item $f_1(X)\equiv X\pmod{X(1-X)}$.
\end{enumerate}
现在设 $(A,\mathfrak m)$ 为完备局部环，$B$ 在 $A$ 上有限，$\varepsilon$ 为 $B/\mathfrak mB$ 中的一个幂等元.关键在于 $B$ 关于 $\mathfrak m$ 进拓扑（由 $x+\mathfrak m^iB$ 生成，其中 $x\in B$、$i\in\mathbb N$）完备（且分离），因为它在 Noether 完备环 $A$ 上有限.
选取 $b\in B$ 使 $\overline b=\varepsilon$，并考虑序列 $(f_n(b))_{n\in\mathbb N}$.因为 $b(1-b)\in\mathfrak mB$，第 ii) 项告诉我们，这是关于 $\mathfrak m$ 进拓扑的一个 Cauchy 列.所以它在 $B$ 中有极限 $\widetilde\varepsilon$.由第 i) 项取极限，得到 $\widetilde\varepsilon^2=\widetilde\varepsilon$；第 iii) 项则给出 $\overline{\widetilde\varepsilon}=\overline b=\varepsilon$.
\end{proof}
\begin{remark*}
这种借助幂等元的方法的好处在于，以上证明在 $B$ 非交换时也适用.例如，它可以证明，对有限群 $G$，$\mathbb F_l[G]$ 的幂等元可以提升到 $\mathbb Z_l[G]$.
\end{remark*}
\subsection{Newton 方法}\label{sec:nt-4-3-3}
设 $(K,|\cdot|)$ 是赋值环为 $\mathcal O$ 的非 Archimedes 完备赋值域.选定首一多项式 $f\in\mathcal O[X]$.当绝对值离散时，Hensel 引理表明 $\overline f$ 的一个孤立根可以提升为 $f$ 的一个根.Newton 方法给出这种根的一种精确而且可算法实现的构造，同时不要求绝对值离散.
\begin{theorem*}
沿用上述记号，设 $x_0\in\mathcal O$ 满足 $|f(x_0)|<|f'(x_0)|^2$.则存在唯一的 $x\in\mathcal O$，使 $f(x)=0$ 且
\begin{align*}
|x-x_0|\le\frac{|f(x_0)|}{|f'(x_0)|}.
\end{align*}
\end{theorem*}
证明与 $\mathbb R$ 上通常的情形类似.对于存在性，令 $\delta:=\frac{|f(x_0)|}{|f'(x_0)|^2}<1$，可以证明递归序列 $x_n:=x_{n-1}-\frac{f(x_{n-1})}{f'(x_{n-1})}$ 对每个 $n$ 满足：
\begin{enumerate}[label=\roman*)]
\item $|f(x_n)|\le\delta|f(x_{n-1})|$.
\item $|f'(x_n)|=|f'(x_{n-1})|$（所以该序列确有定义）.
\end{enumerate}
因此它是一个 Cauchy 列，且其极限 $x$ 具有所要求的性质.
